\section{The invariant derivative of the doubled heat solution}

% BEGIN H75L1
\paragraph{H75L1. The symbol after the doubled substitution and the Gaussian.}
\label{claim:H75L1}
Fix the source frames $x_1,x_2$ in their domain, the ordering
$\mathrm{idx}$, and put $Z_*=\mathrm{dblZ}(x_1,x_2)$.
Let $(P_j)$ satisfy the original column-invariance condition
$\mathrm{InvFam}$. Write $F$ for any holomorphic germ in the two
doubled elliptic blocks $(\zeta,\eta)$ at zero. There is a
constant-coefficient symbol $R$ such that
\[
 \mathcal D_P\{k_0(x_1,x_2;z,w)\,
      F(\mathrm{dblz}(z,w),\mathrm{dblz}'(z,w))\}
  = [R(\partial_\zeta,\partial_\eta)F](0,0).
\]
At each $\tau$ its two $n$-by-$r$ symbol matrices $X_\tau,Y_\tau$
transform by $X_\tau\mapsto\bar u_\tau X_\tau$,
$Y_\tau\mapsto u_\tau Y_\tau$, independently for
$u_\tau\in U(n)$. The polynomial $R$ is invariant under this
action.

\begin{proof}
Identify the $g$ coordinates with $(\tau,a)$ and use zero-based
indices. For $p=X.p_0$, the original linear maps to the doubled
elliptic variables operate on each row $k$ separately. At $\tau$
let $G_\tau$ be the embedded Hermitian matrix and let
$A_{\rm lin},A_{\rm anti}$ be the two source mixed-frame matrices.
The $\zeta$ block takes its input in coordinate $a$ from $z$ when
$a<p$ and from $w$ otherwise; the $\eta$ block takes it from $w$
when $a<p$ and from $z$ otherwise. The respective matrices
multiplying these input columns are
\[
  -\frac{\tau\kappa_0}{2}G_\tau A_{\rm anti}^{-1},
  \qquad
   \frac{\tau\kappa_0}{2}G_\tau^{\mathsf T} A_{\rm lin}^{-1}.
\]
They do not mix rows $k$. Thus a $z$-letter in the second
signature block transforms by $u_\tau$ on its symbol columns, and
a $z$-letter in the first block transforms by $\bar u_\tau$.
The $w$-letters have the two opposite transformations. This
identifies the two substitutions, including the original reversal
of the letter order within its first block.

Let $S(X,Y)$ be the sum of the products obtained by these linear
substitutions in $P_j$ for the $z$-letters and the
coefficient conjugate $\bar P_j$ for the $w$-letters.
By the chain rule this is the symbol on the doubled variables
before the Gaussian multiplier. It is invariant: the defining
identity of $\mathrm{InvFam}$ applies to the two substituted
letter lists, taking the second argument there to be the complex
conjugate of the $w$-list. This turns its conjugated evaluation
into evaluation of $\bar P_j$ on independent $w$-letters. If
symbols are written as columns the substitution in the source
column-action convention uses the adjoint unitary, which ranges
over the same entire group.

At these fixed frames the prefactor, as a function of doubled
variables $v=(\zeta,\eta)$, is
\[
 A(v)=\prod_\tau c_\tau\,
   \exp\!\left(\pi\sum_{k,a,b}
    \eta_{\tau ka}[(\operatorname{Im}Z_{*,\tau})^{-1}]_{ab}
                      \zeta_{\tau kb}\right).
\]
Here $c_\tau=((\det A_{\rm lin}\det A_{\rm anti})^{n/2})^{-1}$
is the original determinant factor. The quadratic pairs a row
column with its dual, so $A$ is invariant under the dual action
on the variables $v$. Absorb it by using formal frequency
variables $\xi=(X,Y)$ and the ordinary bilinear evaluation
$\langle v,\xi\rangle=\sum v_i\xi_i$. Set
\[
 R(\xi)=
 [S(\partial_v)\{A(v)e^{\langle v,\xi\rangle}\}]_{v=0}.
\]
This is a polynomial of degree at most $\deg S$; for $S=0$ take
$R=0$. The Leibniz formula, or coefficient comparison of the
exponential in $\xi$, gives
$[S(\partial)(AF)](0)=[R(\partial)F](0)$ for every germ $F$.
Changing variables contragrediently preserves both the exponent's
bilinear pairing and $A$; $S$ is invariant, so $R$ has the stated
invariance. This includes all contractions from derivatives hitting
the Gaussian.
\end{proof}
% END H75L1

% BEGIN H75L2
\paragraph{H75L2. The invariant contractions are period derivatives.}
\label{claim:H75L2}
Assume $r\le n$. For every holomorphic $H(Z,\zeta,\eta)$
satisfying the native $\mathrm{IsHeat}$ equations and
$H(Z,0,0)=0$ at every $Z$ in the domain, one has
\[
 \mathcal D_P\{k_0(x_1,x_2;z,w)\,
  H(Z_*,\mathrm{dblz}(z,w),\mathrm{dblz}'(z,w))\}=0 .
\]

\begin{proof}
First observe that an invariant polynomial in the symbol variables
of H75L1 is a polynomial in the pairings
$C_{\tau,ab}=\sum_k X_{\tau ka}Y_{\tau kb}$. Here is an
elementary proof for exactly the required range. Its unitary
invariance extends to the complex action
$(X,Y)\mapsto(g^{-\mathsf T}X,gY)$ at each place. For fixed
matrices the difference of the two evaluations is holomorphic in
$g\in\mathrm{GL}_n(\mathbb C)$. For $g=u\exp(tB)$ with
$B$ skew-Hermitian and $u$ unitary it vanishes for real $t$,
hence for complex $t$. Polar decomposition realizes every
invertible matrix this way at $t=i$. Use the same parameter
with independent $B,u$ at the finitely many places.

On the nonempty dense open set where all $C_\tau$ are invertible,
$Y_\tau$ has rank $r$. Carry it to
$Y_0=\binom{1_r}{0}$ by an invertible change of basis. Then
$X_\tau=\binom{C_\tau^{\mathsf T}}{D}$, since the pairings
are unchanged. An element of the stabilizer of $Y_0$ of the form
$\left(\begin{smallmatrix}1&T\\0&1\end{smallmatrix}\right)$
changes the lower block of $X_\tau$ to
$D-T^{\mathsf T}C_\tau^{\mathsf T}$. Choose $T$ to make it zero.
Thus the value of $R$ is the polynomial $s((C_\tau)_\tau)$
obtained by restriction to these explicit representative matrices.
Density makes the identity polynomial everywhere. This works also
for $n=r$, when the lower block is empty, and for $r=0$ the
symbol has only its constant term.

In the source convention the heat equation reads
\[
 B_{\tau,ab}H=
 \sum_k\partial_{\zeta_{\tau ka}}\partial_{\eta_{\tau kb}}H
 =2\pi i\,\partial_{Z_{\tau ab}}H.
\]
All these are constant-coefficient operators on an open complex
domain. Iterating and commuting the derivatives identifies every
monomial in the $B_{\tau,ab}$ on $H$ with the same monomial in
$2\pi i\,\partial_{Z_{\tau ab}}$ on $H$. Restriction to
$(\zeta,\eta)=(0,0)$ commutes with these period derivatives.
The function $Z\mapsto H(Z,0,0)$ vanishes identically, including
its zeroth derivative. Consequently
$[s(B)H](Z_*,0,0)=0$. In this argument the coefficients of $s$
are fixed at the original frame pair: the period derivatives act
only on $H$. This avoids any differentiability-in-$Z$ assumption
on $(\operatorname{Im}Z_*)^{-1}$. Substitution in H75L1 proves the
claim.

For the actual weighted difference in O74E4, O74E5 supplies
$\mathrm{IsHeat}$. The undifferentiated definite Siegel--Weil
identity $H(Z,0,0)=0$ for every $Z$ would now give the required
differentiated comparison. Its proof remains a distinct input,
as do the native Weil and level-system hypotheses of the orbit
calculation. C75E4 derives its absolute convergence at zero from
the finite Weil covariance and unitarity.
\end{proof}
% END H75L2
